\section{Introduction}
\label{sec:intro}
% backgroud
Recent breakthroughs in video generation models~\cite{sora,runway,cogvideox,wan,hunyuanvideo} have shown remarkable progress. Due to their powerful generative capabilities developed through training on large-scale real-world datasets, these models are considered to have the potential to become world models capable of modeling reality~\cite{sora,video-new-language,unisim,worldsimbench}. 
% Among various research directions in this field, interactive long video generation has emerged as a crucial one for two main reasons: (1) Many applications, such as gaming~\cite{genie2,gamengine,gamefactory} and simulation~\cite{gaia-1,vista,gaia-2}, require interactive long video generation, where the videos are generated streamingly controlled by user interactions; (2) By learning from long videos, generative models may acquire deeper knowledge about the world and potentially demonstrate emergent intelligence~\cite{igv-position,igv-survey,lct} \xihui{how this reason is connected to interactive video generation?}. 
% Recent advances in autoregressive video generation~\cite{cdf,dfot,videopoet,emu3,nova,far,framepack} have significantly facilitated research in this field 
% \xihui{How autoregressive video generation is related to interactive long video generation?}.
Among various research directions in this field, interactive long video generation has emerged as a crucial one since many applications, such as gaming~\cite{genie2,gamengine,gamefactory} and simulation~\cite{gaia-1,vista,gaia-2}, require interactive long video generation, where the videos are generated in a streaming manner controlled by user interactions.
Recent works on long video generation~\cite{cdf,dfot,videopoet,emu3,nova,far,framepack} have significantly facilitated research in this field.

% raise the question: memory ability in today's long video gen
Despite these advances, current approaches still face significant challenges in terms of memory capabilities~\cite{igv-position,igv-survey}, which refers to a model's ability to maintain content consistency during continuous video generation, such as preserving the scene when the camera returns to a previously viewed location. Take Oasis~\cite{oasis} as an example: while it can generate lengthy Minecraft gameplay videos, even simple operations like turning left and then immediately right result in completely different scenes. This issue is prevalent across various state-of-the-art methods~\cite{gamengine,wham,dfot,gamefactory}, suggesting that while current approaches can generate videos of extended duration, they struggle to maintain coherent long-term memory of scene content and spatial relationships.

% analyze the question and provide the solution: context as memory
However, in our view, these methods' limitations in memory capabilities are not surprising. This is because when generating each new video frame, these methods can only predict based on a limited number of previous frames. For instance, Diffusion Forcing~\cite{cdf, dfot} can only utilize context from a fixed window of several dozen frames. While this setup works for video continuation, it fails to maintain long-term consistency. 
In the case of video generation, if each frame to be generated could reference all previously generated frames, the generative model could actively select and replicate relevant content from historical frames into the current frame being generated, thus it would be possible to maintain scene consistency in long videos. In other words, \textbf{all previously generated context frames serve as the memory.}

% analyze the proposed solution and analyze its questions
However, the idea of "all historical context as memory" seems intuitive but is impractical for three main reasons:
(1) Including all historical frames in computation would be extremely resource-intensive. 
(2) Processing all historical frames is computationally wasteful since only a small fraction is relevant to the current frame generation.
(3) Processing irrelevant historical frames adds noise that may hinder rather than help current frame generation.
Therefore, a reasonable approach is to retrieve a small number of relevant frames from historical context as conditions for current generation, which we call "\textbf{Memory Retrieval}".

In this work, we propose \textbf{Context-as-Memory} as a solution for scene-consistent interactive long video generation, which includes two simple yet effective designs: (1) Storage format: directly store generated context frames as memory, requiring no post-processing such as feature embedding extraction or 3D reconstruction; (2) Conditioning method: directly incorporate as part of the input through concatenation for context learning, without requiring additional control modules like external adapters or cross attention.
To effectively reduce unnecessary computational overhead and only condition on truly relevant context, we propose \textbf{Memory Retrieval}. Specifically, we introduce a rule-based approach based on camera trajectories. With a camera-controlled video generation model, we can annotate all context frames with camera information based on user's camera control. We can determine co-visibility by checking the FOV (Field of View) overlap based on camera poses at each timestamp along the trajectory, and then use this co-visibility relationship to decide which relevant frames to retrieve.
% To implement this solution, we collected a new dataset consisting of long videos with precise camera annotations from diverse scenes based on Unreal Engine 5.
To implement this solution, we collected a new scene-consistent memory learning dataset using Unreal Engine 5, featuring long videos with precise camera annotations across diverse scenes and camera trajectories. The same regions are captured across different viewpoints and times, enabling both FOV-based retrieval and long-term consistency supervision.

% Provide the method to implement Adaptive Context Compression
% To implement ACC, the key is determining co-visibility relationships between frames. In this work, we introduce a rule-based approach based on camera trajectories. If we know the camera trajectories corresponding to all frames, we can determine co-visibility by checking the FOV overlap based on camera poses at each timestamp along the trajectory. Although this approach doesn't consider occlusion relationships and may not be perfectly precise, it effectively reduces the number of candidate frames. However, camera trajectories are typically unknown to general video generation models. Therefore, to obtain camera trajectories, we introduce camera control mechanisms~\cite{motionctrl,recammaster} into our video generation model, meaning each video segment is generated with user-provided camera trajectories, thus annotating all historical context with camera information. Another benefit of supporting camera control is that we can control the trajectory of new videos to be predicted, allowing users to freely explore the generated virtual world. To implement this solution, we collected long videos with precise camera annotations based on Unreal Engine 5~\cite{ue5}. It's worth emphasizing that ACC implementation doesn't necessarily rely on knowing camera trajectories beforehand. It's possible to develop a self-supervised learnable ACC module, which would make it easier to scale up pre-trained models. This represents a promising direction for future exploration.

Our main contributions can be summarized as follows:
\begin{itemize}
    \item We propose \textbf{Context-as-Memory}, highlighting the direct storage of frames as memory and conditioning via historical context learning for scene-consistent video generation.
    \item To effectively utilize relevant history frames while minimizing costs, we design \textbf{Memory Retrieval}, a specialized rule-based approach using FOV overlap of camera trajectory.
    \item We introduce a long, scene-consistent video dataset with precise camera annotations for memory training, featuring diverse scenes and captions.
    \item Our experiments show superior long video generation memory, significantly outperforming SOTAs and achieving effective memory even in unseen, open-domain scenarios.
\end{itemize}
